\documentclass[aspectratio=169,10pt]{beamer}

\usetheme{Madrid}
\usecolortheme{default}
\setbeamertemplate{navigation symbols}{}
\setbeamertemplate{footline}[frame number]

\usepackage{booktabs}
\usepackage{graphicx}
\usepackage{amsmath}
\usepackage{siunitx}
\usepackage{xcolor}

\graphicspath{{figures/}}

\title[Chemy Thermochemistry]{Chemy Thermochemistry Pipeline}
\subtitle{From reaction evidence to formation thermodynamics}
\author{Chemy project}
\date{\today}

\newcommand{\dH}{\Delta H}
\newcommand{\dG}{\Delta G}
\newcommand{\dHf}{\Delta H_f^\circ}
\newcommand{\dGf}{\Delta G_f^\circ}
\newcommand{\dHr}{\Delta H_r^\circ}
\newcommand{\dGr}{\Delta G_r^\circ}

\begin{document}

\begin{frame}
  \titlepage
\end{frame}

\begin{frame}{Motivation}
  \begin{itemize}
    \item Reaction graphs contain rich thermochemical constraints, but compound-level formation data is sparse.
    \item LLMs can provide broad reaction-level estimates, but the estimates are noisy and locally inconsistent.
    \item Reference datasets such as BURCAT provide accurate anchors for a limited subset of compounds.
    \item Goal: combine noisy reaction evidence and sparse references into a globally consistent formation thermochemistry table.
  \end{itemize}
\end{frame}

\begin{frame}{End-to-End Pipeline}
  \centering
  \includegraphics[width=\textwidth]{pipeline.pdf}
\end{frame}

\begin{frame}{Reaction Data Acquisition}
  \begin{itemize}
    \item Start from a compound catalog with normalized names, CIDs, formulas, structures, and metadata.
    \item Ask LLMs for documented reaction schemes involving target compounds.
    \item Revalidate proposed reactions and keep confidence-scored valid reactions.
    \item Parse names back to CIDs, deduplicate by canonical reaction ID, and balance stoichiometry.
    \item Thermochemistry is estimated only for balanced reactions.
  \end{itemize}
\end{frame}

\begin{frame}{Thermochemistry Evidence}
  \begin{itemize}
    \item For each balanced reaction, sample LLM estimates of \(\Delta H^\circ_{r,298}\) and \(\Delta G^\circ_{r,298}\).
    \item Store per-reaction mean, standard deviation, source model, and uncertainty scale.
    \item Use values in kcal per balanced reaction equation.
  \end{itemize}
  \vspace{0.1em}
  \centering
  \scriptsize
  \resizebox{0.64\textwidth}{!}{\input{tables/counts.tex}}
\end{frame}

\begin{frame}{LLM Reaction Estimate Distribution}
  \centering
  \includegraphics[width=0.92\textwidth]{llm_reaction_distribution.pdf}
\end{frame}

\begin{frame}{LLM Sampling Uncertainty}
  \centering
  \includegraphics[width=0.92\textwidth]{llm_uncertainty.pdf}
\end{frame}

\begin{frame}{BURCAT Reference Extraction}
  \begin{columns}[T,onlytextwidth]
    \column{0.48\textwidth}
      \begin{itemize}
        \item Parse NASA polynomial records from BURCAT.
        \item Read \(\dHf\) from coefficient blocks.
        \item Compute \(\dGf\) using entropy and elemental reference states.
        \item Match species to compounds primarily through CAS numbers.
      \end{itemize}
    \column{0.52\textwidth}
      \includegraphics[width=\textwidth]{burcat_summary.pdf}
  \end{columns}
\end{frame}

\begin{frame}{Global Formation Solve}
  \begin{itemize}
    \item Unknowns are compound formation values \(x_i = \Delta H^\circ_{f,i}\) or \(x_i = \Delta G^\circ_{f,i}\).
    \item Each balanced reaction contributes a linear constraint:
      \[
        \dHr = \sum_i \nu_i \Delta H^\circ_{f,i},
        \qquad
        \dGr = \sum_i \nu_i \Delta G^\circ_{f,i}.
      \]
    \item Element and BURCAT references are additional direct constraints.
    \item Solve a weighted sparse least-squares system, with lower uncertainty for trusted references.
  \end{itemize}
\end{frame}

\begin{frame}{Reference Injection Policies}
  \small
  \begin{itemize}
    \item \textbf{Elements only}: train on elemental zero references; evaluate against all BURCAT compounds.
    \item \textbf{All BURCAT}: train on all available BURCAT references; measures anchor consistency.
    \item \textbf{80/20 split}: train on all elements plus 80\% of BURCAT compounds; evaluate on held-out 20\%.
  \end{itemize}
  \vspace{0.2em}
  \centering
  \includegraphics[width=0.72\textwidth]{policy_mae.pdf}
\end{frame}

\begin{frame}{Policy Metrics}
  \centering
  \resizebox{0.86\textwidth}{!}{\input{tables/metrics.tex}}
  \vspace{0.4em}

  \footnotesize
  Errors are kcal/mol versus BURCAT references. The 80/20 test row is the primary held-out benchmark.
\end{frame}

\begin{frame}{Held-Out Formation Enthalpy}
  \centering
  \includegraphics[width=0.54\textwidth]{dhf_scatter_split_test.pdf}
\end{frame}

\begin{frame}{Held-Out Formation Free Energy}
  \centering
  \includegraphics[width=0.54\textwidth]{dgf_scatter_split_test.pdf}
\end{frame}

\begin{frame}{Error Distributions}
  \centering
  \includegraphics[width=0.86\textwidth]{error_distributions.pdf}
\end{frame}

\begin{frame}{Reaction-Level Correction}
  \small
  \begin{itemize}
    \item The solved formation table induces reaction thermochemistry for every reaction in the graph.
    \item Residuals show how much the global consistency solve moves individual LLM reaction estimates.
  \end{itemize}
  \centering
  \includegraphics[width=0.74\textwidth]{reaction_residuals.pdf}
\end{frame}

\begin{frame}{Example Held-Out Compounds}
  \centering
  \resizebox{0.96\textwidth}{!}{\input{tables/examples.tex}}
\end{frame}

\begin{frame}{Limitations}
  \begin{itemize}
    \item LLM reaction thermochemistry is broad but noisy.
    \item Phase and standard-state assumptions are still implicit in many reactions.
    \item CAS matching gives useful coverage but does not resolve every BURCAT species.
    \item Random holdout can still leak chemically similar compounds through the reaction network.
  \end{itemize}
\end{frame}

\begin{frame}{Next Steps}
  \begin{itemize}
    \item Add stricter reference policies: formula-family split and element-set split.
    \item Improve phase-aware prompting and reference matching.
    \item Calibrate LLM uncertainty using held-out residuals.
    \item Use solved thermochemistry in reaction ranking and UI explanations.
  \end{itemize}
\end{frame}

\begin{frame}{Conclusion}
  \begin{itemize}
    \item Built a modular thermochemistry pipeline over the Chemy reaction graph.
    \item Combined LLM reaction estimates, elemental anchors, and BURCAT references in one sparse solve.
    \item Added policy-based experiments to separate calibration from held-out validation.
    \item Produced reproducible figures and tables for presentation and future benchmarking.
  \end{itemize}
\end{frame}

\end{document}
